% PAGES: 107-108

From (3.9), it follows that a derivative security is replicable if and
only if the price vector $s_\tau$ of the derivative security at time
$\tau$ belongs to the space generated by the price vectors $S_{j\tau}$,
$j = 1:m$, of the $m$ securities at time $\tau$.

Recall from Definition 3.1 that the $m$ securities from a market model
with $m$ securities and $n$ market states at time $\tau > t_0$ are
non--redundant if and only if their price vectors at time $\tau$ are
linearly independent. Thus, the $m$ securities are redundant if and only
if their price vectors are linearly dependent, in which case the payoff
of one of the securities is replicable by using the other $m-1$
securities. Then, from (3.9), it follows that at least one of the $m$
redundant securities can be synthesized using a portfolio containing the
other $m-1$ securities; see section 1.2.1 for an example of a market with
redundant securities.

\section{Arbitrage--free markets}

Arbitrage opportunities arise when there exist investment strategies that
generate guaranteed profit (i.e., with no risk). For a one period market
model with $m$ securities and $n$ market states, such arbitrage
opportunities occur in either one of the two following instances:

(I). If there exists a portfolio with value 0 at time $t_0$ (i.e., with no
set--up cost at time $t_0$), with nonnegative values in every state of the
market at time $\tau$, and with value strictly greater than 0 in at least
one state of the market at time $\tau$.

(II). If there exists a portfolio with negative value at time $t_0$ (i.e.,
with negative set--up cost at time $t_0$, and therefore generating a
positive cash flow at time $t_0$) and with nonnegative values in every
state of the market at time $\tau$.

The following result gives necessary and sufficient conditions for a one
period market model to be arbitrage--free:

\begin{theorem}
A one period market model with $m$ securities and $n$ market states is
arbitrage--free if and only if there exist positive numbers (also called
state prices\footnote{See section 3.4.1 for an explanation of this
nomenclature. In a nutshell, in a complete market, $Q_k$ is the fair price
at time $t_0$ of a bet that state $\omega^k$ will occur at time $\tau$, or
equivalently, the price at time $t_0$ of state $\omega^k$ occurring at
time $\tau$.}) $Q_k > 0$, $k = 1:n$, such that

\begin{equation}
S_{t_0} \;=\; M_\tau Q,
\end{equation}

where $S_{t_0}$ is the price vector of the $m$ securities at time $t_0$,
$M_\tau$ is the payoff matrix at time $\tau$, and $Q = (Q_k)_{k=1:n}$ is
the $n \times 1$ column vector of state prices.
\end{theorem}

Recall that the column form of the payoff matrix $M_\tau$ is $M_\tau =
\mathrm{col}\left(S_\tau^{\omega^k}\right)_{k=1:n}$; cf.\ (3.2). Then,
from (1.7), it follows that (3.10) can be written as

\begin{equation}
S_{t_0} \;=\; \sum_{k=1}^n Q_k S_\tau^{\omega^k}.
\end{equation}

In other words, a one period market model is arbitrage--free if and only
if the price vector of the securities at the initial time $t_0$ is a
linear combination \textit{with positive coefficients} of the price
vectors of the securities in the $n$ possible states at time $\tau$.

Note that, although $Q_k > 0$ for all $k = 1:n$, $Q_k$ is not the
probability for state $\omega^k$ to occur, since $\sum_{k=1}^n Q_k \ne
1$. Thus, formula (3.11) does not state that the price vector of the
securities at the initial time $t_0$ is an expected value of the price
vectors of the securities at time $\tau$.

The proof of Theorem 3.1 involves the Fundamental Theorem of Linear
Programming, which, informally speaking, states that, if a set of linear
inequalities implies another set of linear inequalities, then it does so
trivially, i.e., by linear combinations.

While the proof of Theorem 3.1 is beyond our scope here, we will establish
the sufficiency of condition (3.10) for a non--arbitrage market model to
be arbitrage--free.

% NOTE: this proof is NOT closed on page 108 -- it continues past the
% chunk boundary onto page 109 (PDF index), which is outside this agent's
% page range. The \begin{proof} below is intentionally left open; the
% sibling/next agent covering page 109 must continue this same proof body
% and supply the matching \end{proof} (with no manual qed symbol, per the
% project's proof-environment convention). Do not add a second
% \begin{proof} for this theorem.
\begin{proof}
(Sufficiency of condition (3.10) from Theorem 3.1.) Assume that there
exists a vector $Q = (Q_k)_{k=1:n}$ such that

\begin{equation}
S_{t_0} \;=\; M_\tau Q
\end{equation}

and $Q_k > 0$, for all $k = 1:n$. We will show that neither an arbitrage
of type (I) nor an arbitrage of type (II) can occur.

An arbitrage of type (I) occurs if and only if there exists an $m \times
1$ positions vector $\Theta$ and a state $l$, $1 \le l \le n$, such that

\begin{equation}
V_{t_0} \;=\; \Theta^t S_{t_0} \;=\; 0; \quad
V_\tau^{\omega^k} \;=\; \Theta^t S_\tau^{\omega^k} \;\ge\; 0,\ \forall\,
k=1:n; \quad
V_\tau^{\omega^l} \;=\; \Theta^t S_\tau^{\omega^l} \;>\; 0.
\end{equation}

Recall from (3.2) that $M_\tau = \mathrm{col}\left(S_\tau^{\omega^k}
\right)_{k=1:n}$ and therefore

\begin{equation}
M_\tau Q \;=\; \sum_{k=1}^n Q_k S_\tau^{\omega^k};
\end{equation}

see (1.7). From (3.12) and (3.14), we find that

\begin{equation}
\begin{aligned}
V_{t_0} &= \Theta^t S_{t_0} \;=\; \Theta^t M_\tau Q \;=\; \Theta^t
\left(\sum_{k=1}^n Q_k S_\tau^{\omega^k}\right) \\
&= \sum_{k=1}^n Q_k \Theta^t S_\tau^{\omega^k} \\
&= \sum_{k=1}^n Q_k V_\tau^{\omega^k},
\end{aligned}
\end{equation}

since $V_\tau^{\omega^k} = \Theta^t S_\tau^{\omega^k}$; see (3.13).

Recall from (3.13) that $V_\tau^{\omega^l} > 0$. Also, $Q_l > 0$, and, for
all $1 \le k \le n$, $Q_k > 0$ and $V_\tau^{\omega^k} \ge 0$. Then, it
follows from (3.15) that

\[
V_{t_0} \;=\; \sum_{k=1}^n Q_k V_\tau^{\omega^k} \;\ge\; Q_l
V_\tau^{\omega^l} \;>\; 0,
\]

since $Q_l > 0$, which contradicts the assumption $V_{t_0} = 0$ from the
no--arbitrage condition (3.13).
% Continues on the next page (PDF 109, outside this chunk): the argument
% goes on to rule out an arbitrage of type (II) before \end{proof}.
